\chapter[General Functional Framework]{General Functional Framework}
\section{Introduction}
The previous chapter identified the log-likelihood field
$a(P,Q)=\log\frac{P}{Q}$
as the common geometric object underlying several classical divergences.
This chapter develops a general functional framework in which these
quantities appear as particular evaluations of the same field.
The purpose of this chapter is not to classify all statistical
divergences, but to define a common mathematical language for
functionals acting on the log-likelihood field.
\section{General Definition}
Let $\mu=(\mu_{1},\ldots,\mu_{n})$ be any probability measure on the index set, and let
\[ \Phi:\mathbb{R}^{n}\rightarrow\mathbb{R}
\]
be a real-valued functional.
\begin{definition}
\label{ch:def:28-1}
The functional associated with the pair (\(\mu,\Phi\)) is
\[ \boxed{
F_{\mu,\Phi}(P,Q)
=
E_\mu[\Phi(a)],
} \]
where
$a=\log(P/Q).$
When \(\Phi\) acts componentwise,
\[ F_{\mu,\Phi} = \sum_{i}
\mu_{i},\Phi(a_i). \]
This definition separates two independent choices:
\begin{itemize}
  \item the reference measure \(\mu\);
  \item the observable \(\Phi\).
\end{itemize}
\end{definition}
\section{Classical Examples}
The framework immediately recovers several familiar quantities.
\[
\begin{array}{lll}
\text{Quantity} & \text{Measure} & \text{Functional} \\
\hline
D_{KL}(P\|Q) & P & x \\
D_{KL}(Q\|P) & Q & -x \\
\text{UOD} & \text{Uniform} & (x-E_U[x])^2
\end{array}
\]
The previous chapter established each of these identities rigorously.
\section{The Role of the Reference Measure}
The measure determines how the logarithmic field is sampled.
Typical choices include
\begin{itemize}
  \item the empirical distribution (P);
  \item the reference distribution (Q);
  \item the uniform distribution.
\end{itemize}
Changing the measure changes the interpretation of the same geometric
field without modifying the field itself.
\section{The Role of the Functional}
The functional determines which feature of the field is extracted.
Examples include
\begin{itemize}
  \item expectation;
  \item centered quadratic energy;
  \item higher-order moments;
  \item nonlinear transformations.
\end{itemize}
Thus the measure and the functional represent two independent design
choices.
\section{Relation with f-Divergences}
An f-divergence has the form
\[ D_f(P,Q) = \sum_{i} Q_i
f!\left(\frac{P_i}{Q_i}\right). \]
Since
$\frac{P_i}{Q_i}=e^{a_i},$
one obtains
$D_f(P,Q) = E_Q[f(e^{a})].$
Therefore every f-divergence can be interpreted as a particular
functional of the same log-likelihood field.
This observation does not identify the present framework with the theory
of f-divergences; rather, it shows that both constructions naturally act
on the same underlying logarithmic object.
\section{Geometric Interpretation}
The logarithmic quotient embedding introduced in Part III produces the
field
$a=L(P)-L(Q).$
The present framework interprets statistical quantities as observables
defined on this field.
Accordingly,
\begin{itemize}
  \item geometry generates the field;
  \item functionals generate numerical invariants.
\end{itemize}
The Universal Optimality Defect corresponds to the uniform quadratic
energy of the field, while the KL divergence corresponds to its
expectation under the distribution (P).
\section{Outlook}
The framework introduced here naturally suggests further questions.
Examples include
\begin{itemize}
  \item higher-order moments of the field;
  \item extremal properties under prescribed constraints;
  \item characterization of optimal functionals;
  \item relationships with information geometry and statistical mechanics.
\end{itemize}
These topics are explored in the following chapter.
\section{Summary}
This chapter introduced a general functional framework based on the
log-likelihood field.
Its essential ingredients are
\begin{itemize}
  \item a geometric field \(a=\log(P/Q)\);
  \item a reference measure \(\mu\);
  \item a functional \(\Phi\).
\end{itemize}
Different choices recover different statistical quantities while
preserving a common geometric origin.
